\begin{titlepage}
   \begin{center}
       
       \huge
       \textbf{Firmware Block Implementation for the Proto-DUNE II Hit Finding Architecture with the HLS Tool}
       \vspace{0.6cm}
       
       \Large
       \textbf{Project report by J. Horswill}
       
       \vspace{0.5cm}
       
       \large 
       Supervised by K. Manolopoulos
            
       \vspace{0.5cm}
       
       \large
       Department of Physics and Astronomy\\
       University of Southampton\\
       United Kingdom\\
       \vspace{0.5cm}
       \today
       \vspace{0.5cm}
            
   \end{center}

\begin{abstract}

	The Deep Underground Neutrino Experiment (DUNE) is a large-scale international particle physics project based in North America. It will consist of two underground neutrino detectors placed within the path of an intense neutrino beam fired from Fermilab. The far detector will be 1300 kilometres from the source, in the Sanford Underground Research Facility (SURF). Objectives of the experiment include the study of certain properties of charge-parity violation (CPV) in the lepton sector, as well as the analysis of proton decay, and the investigation into black hole formation through supernova neutrino detection\cite{deep_underground_neutrino_experiment_2020}. This report focusses on the hit finder architecture of the `Trigger Primitive Generation' (TPG) cores within the front-end electronics of ProtoDUNE II. The project objective was to replace the existing firmware algorithms written for the detector's `Field-Programmable Gate Array' (FPGA) boards with designs synthesised via Xilinx's Vivado `High Level Synthesis' (HLS) software. C++ code was used to implement the desired behaviour. This enabled a comparison of the performance and utilization reports of the new designs and those written manually in a Hardware Development Language (VHDL). HLS has the potential to make firmware design more accessible to a wider range of scientists. This could lead to an increased pace in both the development of FPGA-based detector algorithms and the acquisition of important data.
\end{abstract}

\end{titlepage}

